\section{Background}
\label{sec:background}

\subsection{Conditional flow matching}
Flow matching \citep{lipman2022flow} fits a time-dependent vector
field $v_\theta(t, x)$ whose integral transports a simple prior
$p_0$ onto a data distribution $p_1$. With the rectified-flow
coupling used by MMAudio, $x_t = t\, x_1 + (1 - t) x_0$ with
$x_0 \sim \mathcal{N}(0,I)$ and target velocity $u_t = x_1 - x_0$.
Training minimizes
\begin{equation}
\label{eq:fm}
  \LFM(\theta) = \E_{t, x_0, x_1, \mathbf{C}}\;
    \bigl\lVert v_\theta(t, \mathbf{C}, x_t) - u_t \bigr\rVert_2^2,
\end{equation}
where $\mathbf{C}$ is the conditioning triple (CLIP, Synchformer,
text). At inference we Euler-integrate $v_\theta$ from $t{=}0$ to
$t{=}1$ with classifier-free guidance \citep{tong2023improving}.
The MMAudio \texttt{small\_16k} model uses 25 Euler steps, CFG
strength 4.5, and $\sigma_{\min} = 0$. The loss in
Eq.\,(\ref{eq:fm}) is defined \emph{per individual sample}; no
expectation over \emph{pairs} appears, which is the key
structural fact behind the argument in
Sec.\,\ref{sec:intro:pointwise}.

\subsection{MMAudio: architecture and pointwise couplings}
\label{sec:background:mmaudio}
We briefly restate the pieces of MMAudio our work directly sits
alongside, because the precision of our claim depends on what is
and is not present in the upstream objective.

\paragraph{Streams.} Audio latents $x \in \R^{250 \times 20}$
(VAE, 31.25\,fps); CLIP visual tokens $z^{\text{clip}} \in
\R^{64 \times 1024}$ (8\,fps); Synchformer tokens $z^{\text{sync}}
\in \R^{192 \times 768}$ (24\,fps); CLIP-text tokens
$z^{\text{text}} \in \R^{77 \times 1024}$.

\paragraph{MMDiT joint blocks.}
The first $N_1{=}4$ blocks concatenate the (Q, K, V) triples of
audio, CLIP, and text streams and compute one joint scaled
dot-product attention over all tokens. Each modality then reads
off its own slice of the output:
\[
  [q_a; q_v; q_t],
  [k_a; k_v; k_t],
  [v_a; v_v; v_t]
  \;\longrightarrow\;
  \mathrm{softmax}(QK^\top/\sqrt{d}) V.
\]
Every cross-modal inner product is between a single audio token
and a single visual (or text) token --- an \emph{instance-level}
coupling. There are $N_2{=}8$ audio-only fused blocks after the
joint blocks.

\paragraph{Conditional Synchronization Module.}
Synchformer features are upsampled by nearest-neighbor from
24\,fps to the 31.25\,fps latent rate, summed with the broadcast
global condition $c_g$, and injected through a frame-aligned
adaLN with per-token scale/shift. This is still a pointwise
mechanism: audio token $i$ is modulated by video feature at
temporal position $i$. Removing the module costs MMAudio about
50\% of its DeSync margin.

\paragraph{Joint training and empty tokens.}
Audio--visual (VGGSound) and audio--text (AudioCaps, Clotho,
WavCaps) batches are sampled at roughly $1{:}1$. For audio-only
samples, visual streams are replaced by learnable
$\varnothing_v, \varnothing_{\mathrm{syn}}$ tokens. Missing-modality
masking is therefore a \emph{drop-out}, not a distributional
match: it says ``pretend video was blank,'' not ``match this
audio's geometry to the video-free data's geometry.''

\paragraph{Loss.} The training loss is precisely
Eq.\,(\ref{eq:fm}). No contrastive, InfoNCE, alignment, OT,
mutual-information, or reconstruction auxiliary term is used
\citep{cheng2024mmaudio}. This is the gap our GW regularizer
targets.

\subsection{Pointwise alignment vs relational geometry}
\label{sec:background:pointwise-vs-relational}
\citet{wang2020understanding} decompose contrastive objectives
into two forces: an \emph{alignment} force pulling positive
pairs together and a \emph{uniformity} force spreading the
marginal distribution over the hypersphere. Their analysis
makes two facts explicit: (i) an objective built from
independent pairs --- which includes every term in
Eq.\,(\ref{eq:fm}) and every cross-modal mechanism in
MMAudio --- controls the expected
\emph{per-pair} distance but not the joint distribution of
pairwise distances; (ii) matching marginal geometries requires an
objective that sees \emph{multiple} samples simultaneously.
Gromov--Wasserstein is one canonical such objective: by
construction it takes \emph{all} intra-domain pairs as input.

\subsection{Optimal transport and its geometric cousins}
Given two distributions $\mu, \nu$ and a ground cost $c(x,y)$,
the Kantorovich OT problem seeks a coupling $T \in \Pi(\mu,\nu)$
minimizing $\langle C, T \rangle$.
\citet{cuturi2013sinkhorn} added an entropic penalty
$\varepsilon H(T)$, yielding a strictly convex problem solvable
in $O(n^2)$ per iteration by Sinkhorn scaling --- the
computational backbone we reuse inside GW.

\paragraph{Gromov--Wasserstein.} Classical OT requires a single
ground space. The Gromov--Wasserstein distance
\citep{memoli2011gw,peyre2016gromov,peyre2019computational}
generalizes to incommensurable domains by matching
\emph{intra-domain} pairwise distances:
\begin{equation}
\label{eq:gw}
  \mathrm{GW}^2(\DV, \DA, p, q) =
    \min_{T \in \Pi(p,q)}
    \sum_{i,j,k,\ell}
    \bigl( \DV[i,k] - \DA[j,\ell] \bigr)^{2}
    T_{ij}\, T_{k\ell}.
\end{equation}
The objective is quadratic in $T$. \citet{peyre2016gromov}
showed that the entropic relaxation can be solved by mirror
descent: linearize around the current $T$ to get a tangent
linear OT problem with cost $G = -2\, \DV T \DA$, solve with a
few Sinkhorn iterations, repeat. This is exactly what
\texttt{gw\_regularization.entropic\_gw\_loss} implements.

\paragraph{Fused Gromov--Wasserstein.}
\citet{vayer2020fused} introduced FGW, linearly blending GW with
a pointwise cross-domain cost $C_{\text{cross}}$:
\begin{equation}
\label{eq:fgw}
  \LFGW(T) = \alpha \cdot \mathrm{GW}(\DV, \DA; T)
    + (1-\alpha) \cdot \langle C_{\text{cross}}, T \rangle.
\end{equation}
$\alpha = 1$ recovers GW; $\alpha = 0$ recovers Wasserstein on
$C_{\text{cross}}$. FGW gives a single coupling that respects
both relational and pointwise structure; it is one of our four
variants (Sec.\,\ref{sec:design}).

\subsection{Why GW as the relational term}
Several alternatives could populate the role GW plays in our
objective: CKA / HSIC kernel-alignment, centered-cosine
correlation matrices, or a direct MSE between $\DV$ and $\DA$
under a fixed alignment. GW has three properties that make it
the natural choice here:
\begin{enumerate}[leftmargin=*]
  \item \emph{Permutation-free.} The optimal $T$ is discovered,
    not imposed. We do not have to assume sample $i$ on the
    video side corresponds to sample $i$ on the audio side ---
    even though in our batches they do. A permutation-free
    objective makes the regularizer robust to the mixed
    audio-only / audio-visual batches MMAudio trains on, and
    gives an interpretable coupling $\Tstar$ as a byproduct
    (Sec.\,\ref{sec:experiments:analysis}).
  \item \emph{Distributional.} GW returns a distance between
    metric-measure spaces, not just between matrices. Two batches
    that match up to a relabeling have $\LGW = 0$.
  \item \emph{Computationally tractable at the batch size we
    use.} Entropic GW with inner Sinkhorn runs in
    $O(K_{\text{out}} L B^3)$ time dominated by $B \times B$
    matmuls; at $B = 64$ and $K_{\text{out}} = 5$, $L = 20$ this
    is small next to the transformer forward pass.
\end{enumerate}

\subsection{OT in representation and cross-modal learning}
OT has been used repeatedly as a representation signal.
\citet{xu2020gromov} used GW for graph matching and node
embedding; \citet{chen2020graph} used GW-style objectives for
cross-domain alignment; \citet{courty2017joint} used
joint-distribution OT for domain adaptation. Audio-visual
contrastive work
\citep{morgado2021audiovisual,chen2021localizing,afouras2020self}
matches instances pointwise. Our work is closest in
\emph{technique} to \citet{zhou2025fgwclip}, which used FGW as a
contrastive regularizer in enzyme--reaction representation
learning; the computational skeleton is the same but the
downstream task (generative V2A vs retrieval) and the
representations being compared are very different. To our
knowledge, GW has not previously been added as a batch-level
regularizer to a V2A flow-matching model.

\subsection{Video-to-audio synthesis}
Early V2A systems trained per-class or per-dataset generators
\citep{owens2016visually,owens2018audio}. Recent latent-diffusion
and flow-matching systems --- Diff-Foley
\citep{luo2024difffoley}, Frieren \citep{wang2024frieren},
FoleyCrafter \citep{zhang2024foleycrafter}, V-AURA
\citep{viertola2024vaura}, VATT \citep{liu2024vatt}, V2A-Mapper
\citep{wang2024v2amapper}, ReWaS \citep{jeong2024rewas} --- add
large pretrained visual encoders on top of an audio generator.
MMAudio \citep{cheng2024mmaudio} unified multimodal conditioning
in a single MMDiT stack \citep{esser2024scaling} with rotary
position embeddings \citep{su2024roformer}, Synchformer for
sparse temporal cues \citep{iashin2024synchformer}, and a BigVGAN
vocoder \citep{lee2023bigvgan}, and is the current
FD$_{\text{PaSST}}$ / IB-score / DeSync leader. MMAudio is our
base model.

\subsection{Evaluation for V2A}
We follow MMAudio's evaluation via \texttt{av-benchmark}:
FD$_{\text{PaSST}}$ (Fr\'echet distance in PaSST embedding space,
audio fidelity), Inception Score (diversity and
recognizability), IB-score (ImageBind
\citep{girdhar2023imagebind} cross-modal retrieval), and DeSync
(Synchformer-based misalignment). The four metrics probe
fidelity, diversity, semantic coupling, and timing respectively
and can move in opposite directions --- which is why we report
all four throughout.
